\section{Метод потенциальных функций}

Метод потенциальных функций - метрический классификатор, частный случай метода ближайших соседей. Позволяет с помощью простого алгоритма оценивать вес («важность») объектов обучающей выборки при решении задачи классификации.

В общем виде, алгоритм ближайших соседей есть:
\begin{equation*}
	\displaystyle a(x; X^l) = \arg\max_{y\in Y} \sum\limits_{i=1}^l[y^{(i)}=y]w(i,x), 
\end{equation*}

где $w(i,x)$ — вес, степень близости к объекту $x$ его $i$-го соседа.

Метод потенциальных функций заключается в выборе в качестве веса $w(i,x)$, функции следующего вида:

\begin{equation*}
    \displaystyle w(i,x) = \gamma^{(i)} K\left(\frac{\rho(x, x^{(i)})}{h^{(i)}}\right),
\end{equation*}

где:

\begin{itemize}
	\item $K(r)$ -- ядро, не возрастает и положительно на $[0, 1]$,
	\item $\gamma^{(i)} \ge 0$ -- <<заряд>> объекта $x^{(i)}$, 
	\item $h^{(i)} > 0$ -- параметр, задающий <<ширину потенциала>> объекта $x^{(i)}$. 
\end{itemize}

Основная идея метода была навеяна электростатическим взаимодействием элементарных частиц. Известно, что потенциал («мера воздействия») электрического поля элементарной заряженной частицы в некоторой точке пространства пропорционален отношению заряда частицы ($Q$) к расстоянию до частицы ($r$): $\varphi(r) \sim \frac{Q}{r}$. Из-за этого в качестве $K(r)$ часто выступают функции: $\frac{1}{r}$ или $\frac{1}{r+a}$.

Метод потенциальных функций реализует полную аналогию указанного выше примера. При классификации объект проверяется на близость к объектам из обучающей выборки. Считается, что объекты из обучающей выборки «заряжены» своим классом, а мера «важности» каждого из них при классификации зависит от его «заряда» и расстояния до классифицируемого объекта.

\subsection{Подбор параметров}

Отметим, что параметрическую модель зависимости метода потенциальных функций можно существенно упросить (см. задачу ниже):

\begin{equation*}
	\displaystyle a(x; X^l) = \arg\max_{y\in Y} \sum\limits_{i=1}^l[y_i=y]\gamma_i K\left(\frac{\rho(x, x_i)}{h_i}\right).
\end{equation*}

Из выражения выше следует, что в методе потенциальных функций используются две группы параметров: $\{h_i\}$ и $\{\gamma_i\}$.

«Ширина окна потенциала» $h_i$ выбирается для каждого объекта из эмпирических соображений. «Заряд» $\gamma_i$ объектов выборки можно подобрать, исходя из информации, содержащейся в выборке. Ниже приведен алгоритм, который позволяет «обучать» параметры $\gamma_1, \dots, \gamma_n$, то есть подбирать их значения по обучающей выборке $X^l$.

\subsection*{Алгоритм подбора параметров $\{\gamma_i\}$}

\noindent \textbf{Вход:} обучающая выборка из l пар «объект-ответ» -- $X^l=\left((x_1,y_1), \dots, (x_l,y_l) \right)$.

\noindent \textbf{Выход:} значения параметров $\gamma_i$ для всех $i=\overline{1,l}$.

\noindent \textbf{Описание:}

\begin{enumerate}
	\item  Инициализация: $\gamma_i:=0$ для всех $i=\overline{1,l}$; 
	\item Повторять пункты $3-4$, пока эмпирический риск $Q(a,X^l) > \varepsilon$ (то есть пока процесс не стабилизируется):  
	\item Выбрать очередной объект $x_i$ из выборки $X^l$;
	\item Если $a(x_i) \not= y_i$, то $\gamma_i:=\gamma_i+1$;
	\item Вернуть значения $\gamma_i$ для всех $i=\overline{1,l}$.
\end{enumerate}

\subsection{Преимущества и недостатки}

Преимущества метода потенциальных функций:

\begin{itemize}
	\item Метод прост для понимания и алгоритмической реализации;
	\item Порождает потоковый алгоритм; 
	\item Хранит лишь часть выборки, следовательно, экономит память.
\end{itemize}

\noindent Недостатки метода:

\begin{itemize}
	\item Порождаемый алгоритм медленно сходится;
	\item Параметры $\{\gamma_i\}$ и $\{h_i\}$ настраиваются слишком грубо; 
	\item Значения параметров $\gamma_1,\dots,\gamma_l$ зависят от порядка выбора объектов из выборки $X^l$.
\end{itemize}

\subsection{Задачи для самопроверки}

\subsection*{Задача 1.}
Покажите, что параметрическую модель зависимости метода потенциальных функций можно переписать в следующем виде:

\begin{equation*}
	\displaystyle a(x; X^l) = \arg\max_{y\in Y} \sum\limits_{i=1}^l[y_i=y]\gamma_i K\left(\frac{\rho(x, x_i)}{h_i}\right).
\end{equation*}

\subsection*{Решение.}

Ответ следует из независимости суммы от перестановки слагаемых. $\forall y \in Y$:
\begin{equation*}
	\displaystyle 
	\sum\limits_{i=1}^l[y^{(i)}=y]\gamma^{(i)} K\left(\frac{\rho(x, x^{(i)})}{h^{(i)}}\right)
	= \sum\limits_{i=1}^l[y_i=y]\gamma_i K\left(\frac{\rho(x, x_i)}{h_i}\right).
\end{equation*}


\subsection*{Задача 2.}

Покажите, что метод потенциальных функций можно свести к линейному классификатору в случае $Y=\{-1, +1\}$ (бинарная классификация).

\subsection*{Решение.}

Обозначим $\Gamma_y(x) = \sum\limits_{i=1}^l[y^{(i)}=y]w(i,x)$ - оценка близости объекта $x$ к классу $y$. Тогда получим:
\begin{equation*}
	\displaystyle a(x; X^l) = \arg\max_{y\in Y} \Gamma_y(x) = \text{sign} \left(\Gamma_{+1}(x) - \Gamma_{-1}(x)\right) =  \text{sign} \sum\limits_{i=1}^l \gamma_i y_i K\left(\frac{\rho(x, x_i)}{h_i}\right),
\end{equation*}

что совпадает с линейной моделью классификации 
\begin{equation*}
	\displaystyle a(x; X^l) = \text{sign} \sum\limits_{j=1}^n \gamma_j f_j(x),
\end{equation*}

где

\begin{itemize}
	\item $f_j(x) = y_j K\left(\frac{\rho(x, x_j)}{h_j}\right)$ -- новые признаки объекта x,
	\item $\gamma_j$ -- веса линейного классификатора, 
	\item $n = l$ -- число признаков равно числу объектов обучения.
\end{itemize}

\subsection*{Задача 3.}

Привести пример выборки $(x_1, y_1),\dots,(x_l, y_l)$, такой, что метод потенциальных функций построит модель с эмпирическим риском $ Q(a,X^l) = \frac{1}{l} \sum\limits_{i=1}^{l} [a(x_i) \neq y_i] = \frac{1}{l}$.

\subsection*{Решение.}

Рассмотрим $l$-мерное пространство признаков $X=R^l$. В качестве выборки можно взять набор объектов, геометрически находящихся в вершинах $l$-симплекса. Тогда $\rho(x_i, x_j) = \text{const}$ (для простоты положим $\rho(x_i, x_j) = 1$). Остается лишь определить множество ответов: $y_1=+1, y_2=+1, \dots, y_{l-1}=+1, y_l=-1$. Используя результат задачи 2 (положив $h_i=1$), получим:

\begin{equation*}
	\displaystyle a(x; X^l) =  \text{sign} ( \sum\limits_{i=1}^{l-1} \gamma_i - \gamma_l).
\end{equation*}

 Алгоритм подбора параметров $\{\gamma_i\}$ приведет к следующему набору (уже на второй итерации): $\gamma_1=1, \dots, \gamma_{l-1}=0, \gamma_l=1$. Т.е. $\forall x \in X^l$ 
\begin{equation*}
	\displaystyle a(x; X^l) = 1.
\end{equation*}

Воспользовавшись формулой для эмпирического риска получаем требуемое.










